%%%%%%%%%%%%%%%%%%%%%%%%%%%%%%%%%%%%%%%%%%%%%%%%%%%%%%%%%%%%%%%%%%%%%%%%%%%
% Función: capítulo de ejemplo de trabajos previos de la tesis por compendio.
%
% Uso: modifica este fichero para adaptar el contenido a tu trabajo.
%%%%%%%%%%%%%%%%%%%%%%%%%%%%%%%%%%%%%%%%%%%%%%%%%%%%%%%%%%%%%%%%%%%%%%%%%%%

\chapter{Trabajo previo}
\label{cha:compendium-previous-work}

Una revisión bibliográfica puede distinguir un método concreto de localización visual~\cite{Sanabria2026Landmark}, una propuesta de evaluación funcional~\cite{nieva2024proposal} y una aplicación de detección y localización mediante sensado acústico distribuido~\cite{ramirez2026jstars}. La revisión de métodos de aprendizaje automático para vigilancia de tuberías~\cite{tejedor2017machine} ofrece además un ejemplo de revisión del estado del arte, en lugar de un artículo sobre un método individual. Estos ejemplos ilustran formas de organizar las fuentes, pero no implican que todos los trabajos aborden la misma pregunta de investigación.

\section{Antecedentes científicos}
\subsection{Enfoques fundamentales}

Revisa aquí los antecedentes científicos y técnicos de la tesis. Organiza la discusión alrededor de las preguntas de investigación en lugar de presentar una lista aislada de publicaciones, e identifica las limitaciones que motivan el trabajo doctoral.

\subsection{Desarrollos recientes}

Resume los enfoques más relacionados con la tesis y compara sus hipótesis, fortalezas y limitaciones.

\section{Oportunidad de investigación}

Concluye la revisión identificando las cuestiones no resueltas que abordan los objetivos de esta tesis.

%%% Local Variables:
%%% TeX-master: "../../book"
%%% End:
